\documentclass[11pt,a4paper]{article}
\usepackage[margin=23mm,headheight=14pt,headsep=8.5mm,footskip=10mm]{geometry}
\usepackage{amsmath,amssymb}
\usepackage{assets/tex/cm-bea}
\usepackage{xeCJK}
\setmainfont{Latin Modern Roman}
\setCJKmainfont{FandolSong-Regular.otf}[Path=assets/fonts/fandol/,BoldFont=FandolSong-Bold.otf,ItalicFont=FandolSong-Regular.otf]
\setCJKsansfont{FandolSong-Regular.otf}[Path=assets/fonts/fandol/,BoldFont=FandolSong-Bold.otf]
\xeCJKsetup{PunctStyle=quanjiao,CheckSingle=true}
\usepackage{setspace}\setstretch{1.12}
\setlength{\parindent}{2em}\setlength{\parskip}{3pt}
\usepackage{titlesec}
\titleformat{\section}{\centering\normalfont\large\bfseries}{}{0pt}{}
\titlespacing*{\section}{0pt}{1.5\baselineskip}{0.8\baselineskip}
\usepackage{needspace,enumitem,longtable,booktabs,array,calc}
\providecommand{\tightlist}{\setlength{\itemsep}{3pt}\setlength{\parskip}{0pt}}
\usepackage{fancyhdr}\pagestyle{fancy}\fancyhf{}
\fancyhead[L]{\small 数学手记}\fancyhead[R]{\small 楼红卫 B 组 · 第四十批}
\fancyfoot[C]{\thepage}\renewcommand{\headrulewidth}{0.4pt}
\usepackage[hidelinks,unicode,linktoc=all]{hyperref}\usepackage{bookmark}
\hypersetup{pdftitle={楼红卫数学分析 B275—B281：稠密图像、Hardy 不等式与振荡积分},pdfauthor={Vincent Arqain}}
\renewcommand{\contentsname}{题号目录}\setcounter{tocdepth}{1}\setcounter{secnumdepth}{0}
\widowpenalty=10000\clubpenalty=10000\displaywidowpenalty=10000
\predisplaypenalty=10000\interdisplaylinepenalty=10000
\raggedbottom\emergencystretch=2em
\begin{document}
\begin{center}
{\Large\bfseries 楼红卫《数学分析》B 组习题解答\par}
\vspace{7pt}{\large 第四十批\quad B275—B281\par}
\vspace{5pt}{稠密图像、Hardy 不等式与振荡积分\par}
\end{center}
本批收录习题 10.3.B 第 5—8 题及 10.4.B 第 1—3 题。正文依次给出原题、基础知识、思路来源、完整证明、易错条件与方法总结。数学公式采用 CM bea，中文采用 Fandol 宋体；公式编号与网页一致。
\begingroup\setlength{\parskip}{0pt}\tableofcontents\endgroup

\Needspace{9\baselineskip}
\section{问题 B275}\label{ux95eeux9898-b275}

\textbf{原题号：} 习题 10.3.B 第 5
题。\href{https://zj139.github.io/math-notes/files/problems/楼分析习题-v2.pdf\#page=31}{原题，PDF
第 31 页（印刷第 29 页）}

\textbf{题目}　构造 \([0,1]\) 上的函数 \(f\)，使得

\begin{equation}
\overline{\{(x,f(x)):x\in[0,1]\}}=[0,1]\times[0,1].
\tag{1}
\end{equation}

\Needspace{3\baselineskip}\textbf{基础知识}　有理数的可数性与稠密性；可数拓扑基；集合稠密等价于与每个非空开集相交。

\Needspace{3\baselineskip}\textbf{思路来源}　无需让图像经过正方形的每一点，只需让它进入每个小矩形。有理端点矩形只有可数多个，可以逐个满足要求；每次避开已经使用的有限多个横坐标，就不会破坏函数的单值性。

\Needspace{3\baselineskip}\textbf{证明}　把所有满足 \(0<a<b<1\)、\(0<c<d<1\)
且端点为有理数的开矩形列成

\begin{equation}
R_n=(a_n,b_n)\times(c_n,d_n),\qquad n\ge1.
\tag{2}
\end{equation}

递归选取 \(x_n\in(a_n,b_n)\cap\mathbb Q\)，使
\(x_n\notin\{x_1,\ldots,x_{n-1}\}\)。这是可行的，因为非空开区间含无限多个有理数，每次只排除有限多个。取
\(y_n=(c_n+d_n)/2\)，定义

\begin{equation}
f(x)=
\begin{cases}
y_n,&x=x_n\text{ for some }n,\\
0,&x\notin\{x_n:n\ge1\}.
\end{cases}
\tag{3}
\end{equation}

各 \(x_n\) 互异，故定义无歧义，且 \(0\le f\le1\)。记图像为 \(G_f\)，则
\(G_f\cap R_n\ne\varnothing\) 对每个 \(n\) 成立。

任取 \(z\in[0,1]^2\) 及 \(\varepsilon>0\)，即使 \(z\)
位于边界或顶点，\(B(z,\varepsilon)\cap(0,1)^2\)
仍为非空开集，其中含一个有理端点矩形 \(R_n\)。于是
\((x_n,f(x_n))\in B(z,\varepsilon)\)，从而
\(z\in\overline{G_f}\)。反过来，由 \(G_f\subset[0,1]^2\)
及正方形的闭性，\(\overline{G_f}\subset[0,1]^2\)，所以（1）成立。

\Needspace{3\baselineskip}\textbf{易错条件}　闭包符号不可省略：每条竖线与函数图像只有一个交点，而与正方形通常有整个线段。上述函数甚至可测、几乎处处为零；这与图像稠密并不矛盾，因为``几乎处处''和``稠密''分别是测度与拓扑概念。它处处不连续：连续点附近的函数值会限制在一个小区间，无法同时逼近该竖线上的高度
\(0\) 与 \(1\)。

\Needspace{3\baselineskip}\textbf{方法总结}　将不可数多个逼近要求压缩为可数基上的相交要求，再用递归选点保持单值约束。这一构造同时说明：函数在可数集上的修改，足以彻底改变图像的拓扑性质。

\Needspace{9\baselineskip}
\section{问题 B276}\label{ux95eeux9898-b276}

\textbf{原题号：} 习题 10.3.B 第 6
题。\href{https://zj139.github.io/math-notes/files/problems/楼分析习题-v2.pdf\#page=31}{原题，PDF
第 31 页（印刷第 29 页）}

\textbf{题目}　设 \(p>1\)，\(f\in L^p(0,+\infty)\) 且 \(f\ge0\)。证明
Hardy 不等式

\begin{equation}
\int_0^\infty\left(\frac1x\int_0^x f(t)\,dt\right)^pdx
\le\left(\frac p{p-1}\right)^p\int_0^\infty f(x)^p\,dx,
\tag{4}
\end{equation}

其中 \((p/(p-1))^p\) 为最佳常数，等号成立当且仅当
\(\int_0^\infty f(x)\,dx=0\)。

\Needspace{3\baselineskip}\textbf{基础知识}　Hölder
不等式与等号条件；绝对连续函数的分部积分；单调收敛定理。非负 \(L^p\)
函数的全区间 \(L^1\) 积分允许为无穷，其积分为零等价于几乎处处为零。

\Needspace{3\baselineskip}\textbf{思路来源}　记 \(F(x)=\int_0^x f\)。被积函数为
\(x^{-p}F^p\)，对权 \(x^{-p}\) 求原函数，正好把 \(F^p\) 降为
\(F^{p-1}f\)，随后 Hölder 闭合估计。最佳常数对应的形式等号函数是
\(x^{-1/p}\)；它恰处于 \(L^p\) 的临界边缘，因而应截断来逼近等号。

\Needspace{3\baselineskip}\textbf{证明}　记 \(p'=p/(p-1)\)。先设 \(f\) 有界且支于某个
\([a,b]\subset(0,\infty)\)。两端边界项 \(x^{1-p}F(x)^p\)
均趋零，分部积分与 Hölder 给出

\begin{equation}
\begin{aligned}
I&:=\int_0^\infty x^{-p}F(x)^p\,dx
=p'\int_0^\infty\left(\frac{F(x)}x\right)^{p-1}f(x)\,dx\\
&\le p'I^{(p-1)/p}\|f\|_p.
\end{aligned}
\tag{5}
\end{equation}

这里 \(I<\infty\) 已由紧支及 \(p>1\)
保证，故可以除去其幂，得到（4）。对一般 \(f\)，令
\(f_n=\min(f,n)\mathbf1_{[1/n,n]}\)，则 \(f_n\uparrow f\)
几乎处处，相应原函数也递增。由单调收敛把（4）推广到一般情形，同时得到
\(I<\infty\)，没有预先假设待证积分有限。

为审查等号，需要证明（5）对一般 \(f\) 仍是精确恒等式。零端由 Hölder 有

\begin{equation}
0\le\frac{F(x)^p}{x^{p-1}}\le\int_0^x f(t)^p\,dt\longrightarrow0.
\tag{6}
\end{equation}

无穷端固定 \(A>0\)，当 \(x>A\) 时

\begin{equation}
\frac{F(x)}{x^{1-1/p}}
\le\frac{F(A)}{x^{1-1/p}}+
\left(\int_A^\infty f(t)^p\,dt\right)^{1/p}.
\tag{7}
\end{equation}

先令 \(x\to\infty\)，再令
\(A\to\infty\)，便得无穷端边界项也趋零。在有限区间分部积分后取极限即可得到（5）；其中混合积分绝对收敛，由已经证明的
\(I<\infty\) 与 Hölder 保证。

若 \(f\ne0\) 且（4）取等号，则（5）的 Hölder 必取等号，因而 \(f=cF/x\)
几乎处处。结合 \(\|F/x\|_p=p'\|f\|_p\) 得 \(c=1/p'\)，于是

\begin{equation}
\begin{aligned}
F'(x)&=\frac{F(x)}{p'x}\quad\text{a.e.},\\
F(x)&=C x^{1/p'},\qquad
f(x)=\frac C{p'}x^{-1/p}\quad\text{a.e.}
\end{aligned}
\tag{8}
\end{equation}

中间结论可在任意 \([a,b]\subset(0,\infty)\) 上对绝对连续函数
\(x^{-1/p'}F(x)\) 求导得到，因此不要求 \(f\) 连续。非零的（8）不属于
\(L^p(0,\infty)\)，矛盾。故等号当且仅当 \(f=0\)
几乎处处，恰等价于原题所述非负积分为零。

证明最佳性，取 \(R>1\) 及 \(f_R(x)=x^{-1/p}\mathbf1_{[1,R]}(x)\)。则
\(\|f_R\|_p^p=\log R\)，且在 \(1\le x\le R\) 上

\begin{equation}
\frac1x\int_0^x f_R(t)\,dt
=p'x^{-1/p}(1-x^{-1/p'}).
\tag{9}
\end{equation}

由 \(0\le1-(1-u)^p\le pu\)（\(0\le u\le1\)）得

\begin{equation}
0\le\log R-\int_1^R\frac{(1-x^{-1/p'})^p}{x}\,dx
\le p\int_1^\infty x^{-1-1/p'}dx=pp'.
\tag{10}
\end{equation}

因此（4）左端与 \(\|f_R\|_p^p\) 的比值至少趋于
\((p')^p\)；（4）又给出相同上界，最佳常数即为 \((p')^p\)。

\Needspace{3\baselineskip}\textbf{易错条件}　不能由非零极值函数不存在就断定常数不最佳；这里是截断函数列逼近最佳值。无穷端的边界消失不能只用
\(F(x)\le x^{1-1/p}\|f\|_p\)，因为该估计只能给有界性，必须使用（7）的尾部小量。

\Needspace{3\baselineskip}\textbf{方法总结}　分部积分确定常数，Hölder
产生形式极值函数，临界幂的对数发散解释``不取到但可逼近''。先截断证明估计，再审查边界与等号，可以避免循环论证。

\Needspace{9\baselineskip}
\section{问题 B277}\label{ux95eeux9898-b277}

\textbf{原题号：} 习题 10.3.B 第 7
题。\href{https://zj139.github.io/math-notes/files/problems/楼分析习题-v2.pdf\#page=31}{原题，PDF
第 31 页（印刷第 29 页）}

\textbf{题目}　设 \(p>1\)，\(f\in L^p(0,+\infty)\) 且 \(f\ge0\)，又设
\(r>0\)、\(r\ne1\)。证明：当 \(r>1\) 时，

\begin{equation}
\int_0^\infty x^{-r}\left(\int_0^x f(t)\,dt\right)^pdx
\le\left(\frac p{r-1}\right)^p
\int_0^\infty x^{p-r}f(x)^p\,dx;
\tag{11}
\end{equation}

当 \(0<r<1\) 时，

\begin{equation}
\int_0^\infty x^{-r}\left(\int_x^\infty f(t)\,dt\right)^pdx
\le\left(\frac p{1-r}\right)^p
\int_0^\infty x^{p-r}f(x)^p\,dx.
\tag{12}
\end{equation}

\Needspace{3\baselineskip}\textbf{基础知识}　加权 Hölder
不等式；非负函数的单调收敛；原函数与尾积分的导数分别为
\(f\)、\(-f\)。以下记 \(p'=p/(p-1)\)。

\Needspace{3\baselineskip}\textbf{思路来源}　\(\int x^{-r}dx=x^{1-r}/(1-r)\) 的符号随 \(r-1\)
改变。要让分部积分后得到正项且边界消失，\(r>1\)
应取从零点累积的积分，\(r<1\)
应取从无穷端累积的尾积分。这不是可任意交换的两种写法。

\Needspace{3\baselineskip}\textbf{证明}　所有积分先按非负 Lebesgue 积分理解，允许无穷值；记

\begin{equation}
J=\int_0^\infty x^{p-r}f(x)^p\,dx.
\tag{13}
\end{equation}

先对有界且支于 \([a,b]\subset(0,\infty)\) 的 \(f\) 证明。若 \(r>1\)，令
\(F(x)=\int_0^x f\)、\(I=\int_0^\infty x^{-r}F^p\)。零端 \(F=0\)；无穷端
\(F\) 为常数而 \(x^{1-r}\to0\)。因而

\begin{equation}
\begin{aligned}
I&=\frac p{r-1}\int_0^\infty x^{1-r}F^{p-1}f\,dx\\
&=\frac p{r-1}\int_0^\infty
\bigl(x^{-r}F^p\bigr)^{(p-1)/p}
\bigl(x^{p-r}f^p\bigr)^{1/p}dx\\
&\le\frac p{r-1}I^{(p-1)/p}J^{1/p}.
\end{aligned}
\tag{14}
\end{equation}

此时 \(I,J\) 有限，得（11）。若 \(0<r<1\)，令
\(G(x)=\int_x^\infty f\)。无穷端 \(G=0\)；零端 \(G\) 为常数而
\(x^{1-r}\to0\)。利用 \(G'=-f\) 得

\begin{equation}
\begin{aligned}
I&:=\int_0^\infty x^{-r}G^p\,dx
=\frac p{1-r}\int_0^\infty x^{1-r}G^{p-1}f\,dx\\
&\le\frac p{1-r}I^{(p-1)/p}J^{1/p},
\end{aligned}
\tag{15}
\end{equation}

于是（12）成立。对一般 \(f\ge0\)，仍用
\(f_n=\min(f,n)\mathbf1_{[1/n,n]}\uparrow f\)；原函数、尾积分、加权积分均递增，单调收敛直接给出（11）、（12），包括扩充实数情形。

若要求有限值不等式，应明确假设
\(J<\infty\)。这时两种内部积分确为有限：\(r>1\) 时加权 Hölder 给出

\begin{equation}
\int_0^x f(t)dt\le J^{1/p}
\left(\int_0^x t^{(r-p)/(p-1)}dt\right)^{1/p'}<\infty;
\tag{16}
\end{equation}

\(0<r<1\) 时把（16）的积分区间换为
\((x,\infty)\)，右端也有限；指数条件分别恰是 \(r>1\) 与
\(r<1\)。所以在加权范数有限的自然定义域上，甚至无需另加未加权的 \(L^p\)
条件。

\textbf{常数的进一步核查}　两常数也是最佳的。若 \(r>1\)，置
\(a=(r-1)/p\)、\(f_R=x^{a-1}\mathbf1_{[1,R]}\)；若 \(0<r<1\)，置
\(a=(1-r)/p\)、\(f_R=x^{-a-1}\mathbf1_{[1,R]}\)。这些函数均属于原题的
\(L^p\) 类，且两种情形都有 \(J=\log R\)。只积分
\(1\le x\le R\)，左端分别不小于

\begin{equation}
\begin{aligned}
a^{-p}\int_1^R\frac{(1-x^{-a})^p}{x}\,dx
&=a^{-p}\bigl(\log R+O(1)\bigr),\\
a^{-p}\int_1^R\frac{(1-(x/R)^a)^p}{x}\,dx
&=a^{-p}\bigl(\log R+O(1)\bigr).
\end{aligned}
\tag{17}
\end{equation}

两个误差的绝对值均不超过 \(p/a\)，仍由（10）所用的初等估计得到。除以
\(J\) 并令 \(R\to\infty\)，即得相应最佳常数 \(a^{-p}\)。

\Needspace{3\baselineskip}\textbf{易错条件}　原题的 \(f\in L^p\) 并不自动推出
\(J<\infty\)，也不自动保证尾积分有限。例如
\(f(x)=(1+x)^{-\alpha}\)、\(1/p<\alpha\le1\)，虽属于 \(L^p\)，却有
\(\int_x^\infty f=\infty\)；在 \(0<r<1\) 时它的 \(J\)
也为无穷。故若把原题理解为``所有积分必为有限实数''，这就是反例，合理修订是补充
\(J<\infty\)。按非负积分允许无穷的通常约定，原不等式成立，无需改动题面。\(r=1\)
的边界权不衰减，不能代入上述公式。

\Needspace{3\baselineskip}\textbf{方法总结}　由权的原函数决定积分方向，用单调截断处理端点，再以临界幂检验常数。必须区分``扩充实数不等式成立''和``算子在加权空间上有界''这两个层次。

\Needspace{9\baselineskip}
\section{问题 B278}\label{ux95eeux9898-b278}

\textbf{原题号：} 习题 10.3.B 第 8
题。\href{https://zj139.github.io/math-notes/files/problems/楼分析习题-v2.pdf\#page=31}{原题，PDF
第 31 页（印刷第 29 页）}

\textbf{题目}　设 \(p,q>1\) 为对偶数，\(f\in L^q(0,+\infty)\) 且
\(f\ge0\)。令

\begin{equation}
F(f)(x)=\int_x^\infty\frac{f(t)}t\,dt.
\tag{18}
\end{equation}

利用 Hardy 不等式与对偶关系证明

\begin{equation}
\|F(f)\|_{L^q(0,\infty)}
\le\frac p{p-1}\|f\|_{L^q(0,\infty)}.
\tag{19}
\end{equation}

\Needspace{3\baselineskip}\textbf{基础知识}　\(1/p+1/q=1\)；非负函数的 Tonelli 换序；\(L^p\) 与
\(L^q\) 范数对偶。此处 \(p/(p-1)=q\)。

\Needspace{3\baselineskip}\textbf{思路来源}　把尾积分与一个检验函数配对，再交换三角形区域
\(0<x<t\) 的积分次序，尾积分就变成 Hardy
平均算子。这正是``先算伴随，再读范数''的过程。

\Needspace{3\baselineskip}\textbf{证明}　对每个 \(x>0\)，Hölder 保证定义（18）有限：

\begin{equation}
F(f)(x)\le\|f\|_q
\left(\int_x^\infty t^{-p}dt\right)^{1/p}<\infty.
\tag{20}
\end{equation}

对任意非负 \(g\in L^p\)，由 Tonelli 与 B276，

\begin{equation}
\begin{aligned}
\int_0^\infty F(f)(x)g(x)\,dx
&=\int_0^\infty f(t)\left(\frac1t\int_0^t g(x)dx\right)dt\\
&\le\|f\|_q\left\|\frac1t\int_0^t g(x)dx\right\|_p
\le q\|f\|_q\|g\|_p.
\end{aligned}
\tag{21}
\end{equation}

为避免预先假设 \(F(f)\in L^q\)，令
\(u_n=\min(F(f),n)\mathbf1_{[1/n,n]}\)。若 \(u_n\ne0\)，取

\begin{equation}
g_n=\frac{u_n^{q-1}}{\|u_n\|_q^{q-1}},
\qquad \|g_n\|_p=1,
\qquad \int_0^\infty u_ng_n=\|u_n\|_q.
\tag{22}
\end{equation}

将 \(g_n\) 代入（21），利用 \(u_n\le F(f)\) 得
\(\|u_n\|_q\le q\|f\|_q\)；\(u_n=0\)
时该式显然成立。单调收敛给出（19）。

事实上常数 \(q\) 最佳。取 \(f_R(t)=t^{-1/q}\mathbf1_{[1,R]}\)，则
\(\|f_R\|_q^q=\log R\)，并且

\begin{equation}
F(f_R)(x)=q x^{-1/q}\bigl(1-(x/R)^{1/q}\bigr),
\qquad1\le x\le R.
\tag{23}
\end{equation}

由（17）同样的估计，\(\|F(f_R)\|_q^q\ge q^q(\log R-O(1))\)，因此算子范数不能小于
\(q\)。

\Needspace{3\baselineskip}\textbf{易错条件}　对偶后改变的是函数所在空间，常数仍是 Hardy 算子在
\(L^p\) 上的范数 \(q\)，不能误写为 \(p\)。Tonelli
可以在尚不知道配对有限时使用；随后（21）证明它有限。不能直接把未证属于
\(L^q\) 的 \(F(f)\) 代入范数对偶公式而省略截断。

\Needspace{3\baselineskip}\textbf{方法总结}　积分区域的转置把``向前平均''变为``向后加权尾积分''。先通过非负换序建立配对恒等式，再以截断范数检验函数完成有界性证明。

\Needspace{9\baselineskip}
\section{问题 B279}\label{ux95eeux9898-b279}

\textbf{原题号：} 习题 10.4.B 第 1
题。\href{https://zj139.github.io/math-notes/files/problems/楼分析习题-v2.pdf\#page=32}{原题，PDF
第 32 页（印刷第 30 页）}

\textbf{题目}　证明

\begin{equation}
\int_0^\infty\frac{e^{-x}-\cos x}{x}\,dx=0.
\tag{24}
\end{equation}

\Needspace{3\baselineskip}\textbf{基础知识}　Dirichlet
判别；含参积分求导；指数阻尼与一致尾估计。各项在零点单独积分都发散，必须保留相消结构。

\Needspace{3\baselineskip}\textbf{思路来源}　余弦没有指数衰减，直接求导容易丢失收敛性。先加入
\(e^{-ax}\)，使参数计算在绝对收敛范围内进行，再用与 \(a\)
无关的尾估计去掉阻尼。

\Needspace{3\baselineskip}\textbf{证明}　零点处
\(e^{-x}-\cos x=-x+O(x^2)\)，故被积函数有有限极限；无穷端的 \(e^{-x}/x\)
绝对可积，\(\cos x/x\) 由 Dirichlet 判别收敛。因此（24）为普通反常积分。

对 \(a>0\)、\(b\ge0\) 定义

\begin{equation}
J(a,b)=\int_0^\infty\frac{e^{-x}-e^{-ax}\cos(bx)}x\,dx.
\tag{25}
\end{equation}

该积分绝对收敛。由 \(e^{-x}-e^{-ax}=\int_1^a x e^{-tx}dt\)（\(a<1\)
时取有向积分），在 \(t\) 的紧区间内换序得

\begin{equation}
J(a,0)=\int_1^a\frac{dt}{t}=\log a.
\tag{26}
\end{equation}

固定 \(a>0\)，求导后的被积函数受 \(e^{-ax}\) 控制，故

\begin{equation}
\frac{\partial J}{\partial b}(a,b)
=\int_0^\infty e^{-ax}\sin(bx)dx
=\frac b{a^2+b^2}.
\tag{27}
\end{equation}

最后一个等式由两次分部积分或直接求原函数得到。对 \(b\) 积分，得

\begin{equation}
J(a,b)=\log a+\frac12\log\frac{a^2+b^2}{a^2}
=\frac12\log(a^2+b^2).
\tag{28}
\end{equation}

取 \(b=1\)。对 \(a\ge0\)，权 \(w_a(x)=e^{-ax}/x\)
单调趋零；分部积分给出一致尾估计

\begin{equation}
\left|\int_R^\infty\frac{e^{-ax}\cos x}{x}dx\right|
\le2w_a(R)\le\frac2R.
\tag{29}
\end{equation}

在有限区间上，\(0<a\le1\) 时
\(|e^{-x}-e^{-ax}\cos x|/x\le1+x/2\)，因此可用控制收敛令
\(a\downarrow0\)。再由（29）及指数项的可积尾部去掉截断，得到原积分等于
\(\lim_{a\downarrow0}J(a,1)=0\)。

\Needspace{3\baselineskip}\textbf{易错条件}　式（24）不能写成两个从零开始的独立积分相减；正确操作是在同一截断区间上相减。阻尼积分的极限也不能仅凭点态收敛与原积分收敛而宣称相等，（29）给出了所需的统一控制。

\Needspace{3\baselineskip}\textbf{方法总结}　相消解决零端奇点，指数阻尼解决无穷端求导，统一尾估计负责返回原积分。这三个环节各自承担不同的收敛任务。

\Needspace{9\baselineskip}
\section{问题 B280}\label{ux95eeux9898-b280}

\textbf{原题号：} 习题 10.4.B 第 2
题。\href{https://zj139.github.io/math-notes/files/problems/楼分析习题-v2.pdf\#page=32}{原题，PDF
第 32 页（印刷第 30 页）}

\textbf{题目}　计算

\begin{equation}
I=\int_0^\infty\frac{\cos x-\cos(x^2)}x\,dx.
\tag{30}
\end{equation}

\Needspace{3\baselineskip}\textbf{基础知识}　有限截断上的换元；Euler 常数
\(\gamma=\lim_{n\to\infty}(H_n-\log n)\)，其中
\(H_n=\sum_{k=1}^n1/k\)；对数奇点的可积性。

\Needspace{3\baselineskip}\textbf{思路来源}　换元 \(t=x^2\) 会把下限 \(\varepsilon\) 变成
\(\varepsilon^2\)。两个积分的对数发散由此抵消，剩下的有限常数才是答案。B279
可以把振荡积分的有限部分转移到指数积分，从而引出 Euler 常数。

\Needspace{3\baselineskip}\textbf{解}　零点处分子为
\(-x^2/2+O(x^4)\)；无穷端的两个积分分别收敛，第二个通过 \(t=x^2\) 化为
\(\frac12\int\cos t/t\,dt\)。故（30）作为普通反常积分存在。

先证明所需常数公式

\begin{equation}
\int_\varepsilon^\infty\frac{e^{-t}}t\,dt
=-\log\varepsilon-\gamma+o(1)
\quad(\varepsilon\downarrow0).
\tag{31}
\end{equation}

为使常数来源明确，考虑

\begin{equation}
A_n=\int_0^n(1-t/n)^n\log t\,dt.
\tag{32}
\end{equation}

把被积函数在 \(t>n\) 处补零，绝对值受 \(e^{-t}|\log t|\) 控制，故
\(A_n\to\int_0^\infty e^{-t}\log t\,dt\)。换元 \(t=nu\)，并利用

\begin{equation}
\begin{aligned}
-\int_0^1(1-u)^n\log u\,du
&=\frac1{n+1}\int_0^1\frac{1-(1-v)^{n+1}}v\,dv\\
&=\frac{H_{n+1}}{n+1},
\end{aligned}
\tag{33}
\end{equation}

得

\begin{equation}
A_n=\frac n{n+1}(\log n-H_{n+1})\longrightarrow-\gamma.
\tag{34}
\end{equation}

式（33）先用 \(-\log u=\int_u^1dv/v\)
非负换序，再用有限几何级数积分；因此没有引用未证的特殊函数恒等式。\(\gamma\)
的存在由积分比较可知：\(H_n-\log n\) 单调递减且有下界。

在 \([\varepsilon,R]\) 上分部积分并令 \(R\to\infty\)，有

\begin{equation}
\int_\varepsilon^\infty e^{-t}\log t\,dt
=e^{-\varepsilon}\log\varepsilon+
\int_\varepsilon^\infty\frac{e^{-t}}t\,dt.
\tag{35}
\end{equation}

由（34）、（35）及 \((e^{-\varepsilon}-1)\log\varepsilon\to0\)
得（31）。又由 B279 的总积分为零，

\begin{equation}
\int_\varepsilon^\infty\frac{\cos t}t\,dt
=\int_\varepsilon^\infty\frac{e^{-t}}t\,dt+o(1)
=-\log\varepsilon-\gamma+o(1).
\tag{36}
\end{equation}

这里误差等于
\(\int_0^\varepsilon(e^{-t}-\cos t)/t\,dt\)，确实趋零。对固定
\(\varepsilon>0\) 在第二项换元，再令 \(\varepsilon\downarrow0\)，便得

\begin{equation}
\begin{aligned}
I&=\lim_{\varepsilon\downarrow0}\left[
\int_\varepsilon^\infty\frac{\cos t}t\,dt
-\frac12\int_{\varepsilon^2}^\infty\frac{\cos t}t\,dt\right]\\
&=\lim_{\varepsilon\downarrow0}
\left[-\log\varepsilon-\gamma
-\frac12(-2\log\varepsilon-\gamma)+o(1)\right]
=-\frac\gamma2.
\end{aligned}
\tag{37}
\end{equation}

\Needspace{3\baselineskip}\textbf{易错条件}　不能把换元后的两个下限都写成零，再形式上合并成
\(\frac12\int_0^\infty\cos t/t\,dt\)；这个积分在零点发散。有限部分的常数会随系数保留下来，并不会被对数项的抵消一同消去。

\Needspace{3\baselineskip}\textbf{方法总结}　遇到同一函数的不同幂尺度，保留变换后的真实截断端点。先分离对数发散，再计算有限部分，往往比直接寻找原函数更有效。

\Needspace{9\baselineskip}
\section{问题 B281}\label{ux95eeux9898-b281}

\textbf{原题号：} 习题 10.4.B 第 3
题。\href{https://zj139.github.io/math-notes/files/problems/楼分析习题-v2.pdf\#page=32}{原题，PDF
第 32 页（印刷第 30 页）}

\textbf{题目}　计算

\begin{equation}
\int_{-\infty}^{+\infty}\sin(t^2)\,dt.
\tag{38}
\end{equation}

\Needspace{3\baselineskip}\textbf{基础知识}　Gauss 积分；振荡积分的 Dirichlet
判别；对参数求导及一阶线性微分方程。以下用复数记号合并正弦、余弦，但只使用实参数积分与分部积分。

\Needspace{3\baselineskip}\textbf{思路来源}　相位 \(t^2\) 与 Gauss 核有相同的二次结构。加入
\(e^{-at^2}\) 后，把正弦和余弦合为一个积分，参数求导产生的 \(t^2\)
正好能通过分部积分消去，从而得到可解的微分方程。

\Needspace{3\baselineskip}\textbf{解}　换元 \(u=t^2\) 得

\begin{equation}
\int_R^S\sin(t^2)dt
=\frac12\int_{R^2}^{S^2}\frac{\sin u}{\sqrt u}\,du.
\tag{39}
\end{equation}

由 Dirichlet
判别，正半轴积分收敛；负半轴由偶性同样收敛。所以（38）不是仅以主值意义存在。

对 \(a>0\)、\(b\ge0\)，令

\begin{equation}
K_a(b)=\int_0^\infty e^{-(a-ib)t^2}dt.
\tag{40}
\end{equation}

绝对收敛使参数求导合法；导数受 \(t^2e^{-at^2}\) 控制。积分
\(\frac{d}{dt}(t e^{-(a-ib)t^2})\)，两端边界均为零，故

\begin{equation}
\begin{aligned}
K_a'(b)&=i\int_0^\infty t^2e^{-(a-ib)t^2}dt,\\
K_a(b)&=2(a-ib)\int_0^\infty t^2e^{-(a-ib)t^2}dt,\\
K_a'(b)&=\frac{i}{2(a-ib)}K_a(b).
\end{aligned}
\tag{41}
\end{equation}

初值为 \(K_a(0)=\sqrt\pi/(2\sqrt a)\)：正的 Gauss
积分平方后，在第一象限使用极坐标，得到

\begin{equation}
K_a(0)^2=\int_0^{\pi/2}\int_0^\infty e^{-ar^2}r\,dr\,d\theta
=\frac\pi{4a}.
\tag{42}
\end{equation}

解（41）得到

\begin{equation}
K_a(b)=\frac{\sqrt\pi}{2\sqrt{a-ib}}.
\tag{43}
\end{equation}

平方根取从 \(b=0\) 的正根连续延伸的分支；由于 \(a>0\)，\(a-ib\)
始终位于右半平面，分支无歧义。也可直接对 \(\sqrt{a-ib}\,K_a(b)\)
求导，验证其恒等于 \(\sqrt\pi/2\)。

取 \(b=1\)，令 \(a\downarrow0\)。为证明阻尼极限等于原积分，令
\(w_a(u)=e^{-au}/(2\sqrt u)\)。对 \(a\ge0\)，它单调趋零，分部积分给出

\begin{equation}
\left|\int_R^\infty e^{-at^2}e^{it^2}dt\right|
=\left|\int_{R^2}^\infty w_a(u)e^{iu}du\right|
\le2w_a(R^2)\le\frac1R.
\tag{44}
\end{equation}

结合有限区间上的控制收敛，可以在（43）取极限并取虚部：

\begin{equation}
\begin{aligned}
\int_0^\infty\sin(t^2)dt
&=\operatorname{Im}\left(\frac{\sqrt\pi}{2}e^{i\pi/4}\right)
=\sqrt{\frac\pi8},\\
\int_{-\infty}^{+\infty}\sin(t^2)dt
&=2\sqrt{\frac\pi8}=\sqrt{\frac\pi2}.
\end{aligned}
\tag{45}
\end{equation}

\Needspace{3\baselineskip}\textbf{易错条件}　偶函数的全实轴积分是半轴积分的两倍，不能漏掉这个因子。积分并非绝对收敛：由（39），在
\([k\pi+\pi/6,k\pi+5\pi/6]\) 上有
\(|\sin u|\ge1/2\)，各段绝对积分下界为常数乘
\(k^{-1/2}\)，其和发散。因此（43）的阻尼极限必须由（44）证明，不能用无可积控制函数的全区间控制收敛。

\Needspace{3\baselineskip}\textbf{方法总结}　二次振荡与 Gauss
核相结合，可将积分转化为参数微分方程。复数记号只负责合并两类实积分；真正保证推导成立的是正阻尼下的绝对收敛，以及去阻尼时的一致尾估计。

\Needspace{9\baselineskip}
\section{七题方法总结}\label{ux4e03ux9898ux65b9ux6cd5ux603bux7ed3}

\begin{longtable}[]{@{}
  >{\raggedright\arraybackslash}p{(\columnwidth - 4\tabcolsep) * \real{0.3333}}
  >{\raggedright\arraybackslash}p{(\columnwidth - 4\tabcolsep) * \real{0.3333}}
  >{\raggedright\arraybackslash}p{(\columnwidth - 4\tabcolsep) * \real{0.3333}}@{}}
\toprule\noalign{}
\begin{minipage}[b]{\linewidth}\raggedright
题号
\end{minipage} & \begin{minipage}[b]{\linewidth}\raggedright
核心结果
\end{minipage} & \begin{minipage}[b]{\linewidth}\raggedright
关键方法
\end{minipage} \\
\midrule\noalign{}
\endhead
\bottomrule\noalign{}
\endlastfoot
B275 & 图像闭包为单位正方形 & 可数基、互异横坐标递归选点 \\
B276 & Hardy 最佳常数及仅零函数取等号 & 分部积分、Hölder、临界幂截断 \\
B277 & 两种加权 Hardy 不等式及最佳常数 &
权的原函数决定积分方向，单调截断 \\
B278 & 对偶尾积分算子的范数为对偶指数 & Tonelli
换序、截断对偶检验函数 \\
B279 & 指数与余弦抵消积分为零 & 参数求导、指数阻尼、一致尾估计 \\
B280 & 不同幂尺度的积分为负半个 Euler 常数 &
保留截断端点、计算对数有限部分 \\
B281 & 全实轴 Fresnel 积分为 \(\sqrt{\pi/2}\) & Gauss
初值、参数微分方程、去阻尼 \\
\end{longtable}

\end{document}
